\newpage
\section{Conclusion}\label{sec:conclusion}

The central finding of this thesis is that, in the sample studied, momentum's cross-sectional predictability shifts with market conditions: which lookback horizons carry predictive power, and in which direction, depends on the prevailing market regime. The regime signal captures this shift not by predicting returns directly but by conditioning how the model jointly uses the full momentum term structure -- a multi-dimensional interaction that only a nonlinear model can exploit. These findings complement the exposure-scaling approach of \citet{DanielMoskowitz2016} and \citet{BarrosoSantaClara2015} with a stock-selection channel: rather than adjusting how much to hold, the regime signal adjusts which stocks to hold.

\subsection{Contributions}\label{sec:contributions}

The three contributions outlined in Section~\ref{sec:introduction} are supported by the empirical evidence. The term-structure reorganisation is visible directly in z-score profiles and does not depend on the modelling framework. The context-variable mechanism is confirmed by SHAP analysis showing that $\pi_t^{\text{filter}}$ conditions the joint use of momentum horizons rather than predicting returns directly. The nonlinearity requirement is established by the failure of both classification-based and regression-based linear models on identical features. Together, these findings suggest that the fragility of momentum documented in prior work may partly reflect unconditional models averaging over a regime-dependent structure, rather than an intrinsic deficiency of the anomaly itself.

\subsection{Scope and Generalisability}\label{sec:generalisability}

The strategy's primary vulnerability is a prolonged fundamental deterioration where beaten-down stocks do not recover. The stress test in Section~\ref{sec:stress_test} quantifies this: a 12-month recession reduces the worst-case Sharpe to 0.75 with a drawdown of $-$51\%, while a 24-month depression pushes the worst-case Sharpe to 0.44 with a $-$70\% drawdown. The strategy survives because calm-period continuation provides a floor, but a sustained downturn where reversal consistently fails would be painful. The 2011--2025 test period does not contain such a scenario.

Several limitations bound what can be concluded. The test period (2011--2025, 167 months) does not contain a prolonged recession, though the 2008--2009 out-of-sample test (Appendix~\ref{app:gfc_oos}) confirms that M2 would have survived the momentum crash. Performance is strongest in 2021--2025 (Sharpe 1.69), a period with frequent regime transitions; in calm-dominated years the model earns modest returns from standard continuation. A permanently calm market would reduce the strategy to ordinary momentum.

The Newey--West $t$-statistic for M2's excess return over the market is 1.83 (marginally significant), but this is expected for a market-neutral strategy with near-zero beta. The economically meaningful test is the factor model alpha, which is significant across all specifications ($t > 4$).

The analysis assumes frictionless short selling and uses a uniform 10 basis point transaction cost. The sample is limited to US equities. The causal interpretation that the regime signal conditions momentum's term structure is supported by multiple lines of evidence (SHAP, tree paths, ablation, LR-XGB divergence) but is inferred from observational data and should be interpreted as consistent evidence rather than definitive proof.

\subsection{Future Work}\label{sec:future_work}

Several directions could extend this framework.

\paragraph{Regime signal improvements}
The current signal $\pi_t^{\text{filter}}$ estimates panic probability at time $t$. A predictive variant $P(s_{t+1} = 1 \mid \mathbf{z}_{1:t})$, obtainable from the transition matrix, could reposition the portfolio \emph{before} regime shifts. Separately, the fixed HMM parameters may become miscalibrated if markets undergo structural change; rolling or expanding-window re-estimation would address this. A sequential model that conditions on the \emph{trajectory} of $\pi_t^{\text{filter}}$ -- distinguishing fresh stress from prolonged deterioration -- could improve detection of reversal failure.

\paragraph{Robustness and generalisation}
Testing the framework on international markets (Europe, Japan, emerging markets) would establish whether the term-structure reorganisation is a universal feature of momentum. Comparing XGBoost against alternative nonlinear models (Random Forests, neural networks) would clarify whether the finding is about the data or the model architecture. The HMM's Gaussian emissions could be replaced with Student-$t$ or skew-normal distributions to better accommodate tail events.

\paragraph{Implementation extensions}
The strategy's vulnerability to prolonged fundamental deterioration could be mitigated by an exposure-scaling circuit breaker that reduces overall position size during sustained bear markets. The feature set may itself be regime-dependent: fundamentals might add value during calm if used as a quality filter rather than as a direct model input. Finally, incorporating stock-specific transaction costs would provide a more realistic picture of implementability.

